%
% menu_presence_locks.tex
%
% Companion to docs/menu_presence.tex: the LOCK ORDER the fixed design
% requires, checked for deadlock-freedom.
%
% menu_presence.tex makes every presence-changing registry transition one
% atomic step. That atomicity is only real if the code can actually hold the
% locks it needs without a cycle. The locks in play:
%
%   StoreLock   HubDisplay._lock, a re-entrant write lock. Held by
%               identify_client, register_client, drop_connection, apply,
%               _depart_lapsed, and by bind_departure_sink. The departure
%               sinks run inside it.
%   RegLock     HubClientRegistry._lock, a plain (non-re-entrant) lock. Taken
%               INSIDE StoreLock by every writer (record, attach_listener,
%               discard, reap_lapsed_locked), and ALONE by the readers
%               (named_sessions, live_sessions, session_of).
%   (marker)    The DirtyMarker's own lock, taken by mark_menus. It is a leaf:
%               nothing is acquired while it is held, so it contributes no
%               edge to the acquisition graph and is not modelled.
%
% Re-entrancy of StoreLock contributes no edge either: a thread already
% holding it never waits on it, so a nested acquire is a self-loop. It is
% modelled by omission.
%
% THE FIXED DESIGN: every presence-changing writer enters through HubDisplay
% (StoreLock first, then RegLock); the registry returns "presence changed" and
% the caller, still inside StoreLock, marks and binds. Readers take RegLock
% alone and never wait for the store. The acquisition graph has the single
% edge StoreLock -> RegLock: acyclic, so deadlock-free.
%
% THE REJECTED DESIGN (menu_presence_locks_inverted_buggy.tex): the registry
% itself fires the mark/bind hook from inside record(), i.e. while holding
% RegLock, and binding a sink takes StoreLock. Writers that enter through the
% registry then hold RegLock and want StoreLock, against HubDisplay's writers
% that hold StoreLock and want RegLock: a cycle. ProB finds the deadlock.
%
% Type-check:  fuzz docs/menu_presence_locks.tex
% Model-check: probcli docs/menu_presence_locks.tex -model_check \
%                -p DEFAULT_SETSIZE 2
% Expected: no deadlock, no invariant violation (the fixed order).
%

\documentclass[a4paper,10pt,fleqn]{article}
\usepackage[margin=1in]{geometry}
\usepackage{fuzz}

\begin{document}

\title{lux Menu Presence: the Lock Order}
\maketitle

\begin{zed}
  [T]
\end{zed}

A thread: an MCP tool thread, the listen leg's event-loop thread, the lease
sweep, the replicator tick. Two suffice for a cycle.

\begin{zed}
  PC ::= pidle | pstore | pboth | preg
\end{zed}

A thread's program counter: holding nothing; holding the store lock and
wanting the registry's; holding both; holding the registry lock alone (a
read).

\begin{schema}{Locks}
  pc : T \fun PC \\
  storeHeld : \power T \\
  regHeld : \power T
\where
  \# storeHeld \leq 1 \\
  \# regHeld \leq 1 \\
  storeHeld = \{ t : T | pc~t = pstore \lor pc~t = pboth \} \\
  regHeld = \{ t : T | pc~t = pboth \lor pc~t = preg \}
\end{schema}

The lock tables are derived from the program counters, so they can never
disagree with them; the cardinality bounds are mutual exclusion.

\begin{schema}{Init}
  Locks'
\where
  pc' = T \cross \{ pidle \} \\
  storeHeld' = \emptyset \\
  regHeld' = \emptyset
\end{schema}

\begin{schema}{BeginWrite}
  \Delta Locks \\
  t? : T
\where
  pc~t? = pidle \\
  storeHeld = \emptyset \\
  pc' = pc \oplus \{ t? \mapsto pstore \} \\
  storeHeld' = \{ t? \} \\
  regHeld' = regHeld
\end{schema}

A presence-changing writer takes the store lock first (HubDisplay's
\texttt{identify\_client}, \texttt{\_depart}, \texttt{\_depart\_lapsed}, and
the read loop once it is routed through HubDisplay).

\begin{schema}{TakeRegistry}
  \Delta Locks \\
  t? : T
\where
  pc~t? = pstore \\
  regHeld = \emptyset \\
  pc' = pc \oplus \{ t? \mapsto pboth \} \\
  storeHeld' = storeHeld \\
  regHeld' = \{ t? \}
\end{schema}

\begin{schema}{Finish}
  \Delta Locks \\
  t? : T
\where
  pc~t? = pboth \\
  pc' = pc \oplus \{ t? \mapsto pidle \} \\
  storeHeld' = \emptyset \\
  regHeld' = \emptyset
\end{schema}

The upsert or removal happens under both locks; the registry returns, and the
caller marks and binds under the store lock it still holds, then releases
both. (The marker's leaf lock, taken in between, is not modelled.)

\begin{schema}{BeginRead}
  \Delta Locks \\
  t? : T
\where
  pc~t? = pidle \\
  regHeld = \emptyset \\
  pc' = pc \oplus \{ t? \mapsto preg \} \\
  storeHeld' = storeHeld \\
  regHeld' = \{ t? \}
\end{schema}

\begin{schema}{EndRead}
  \Delta Locks \\
  t? : T
\where
  pc~t? = preg \\
  pc' = pc \oplus \{ t? \mapsto pidle \} \\
  storeHeld' = storeHeld \\
  regHeld' = \emptyset
\end{schema}

A reader (\texttt{named\_sessions}, the replicator's roster read) takes the
registry lock alone and releases it without ever asking for the store.

Deadlock-freedom is the claim that some operation is enabled in every
reachable state: with the order above, a thread holding \texttt{regHeld}
either finishes (\texttt{pboth}) or is a reader that releases
(\texttt{preg}); neither waits on anything, so the thread blocked on
\texttt{regHeld} is always eventually enabled.

\begin{verbatim}
  fuzz docs/menu_presence_locks.tex
  probcli docs/menu_presence_locks.tex -model_check -p DEFAULT_SETSIZE 2
\end{verbatim}

\end{document}
